% !TEX program = pdflatex
\documentclass[11pt]{article}
\usepackage{support-notes}
\title{A rank criterion for a cubic higher product to change bounded profiles}
\author{\workingauthor}\date{\coursedate}
\begin{document}\maketitle
\begin{abstract}
In a matrix family of finite-dimensional commutative local two-term dg
algebras with constant graded cohomology, the formal two-element profile
poset persists on every deficient-rank stratum. Precisely the square
invertible matrices give the full noncommutative localisation poset.
A general obstruction for radical-square-zero cohomology explains why
many nonzero ternary products still yield only two profiles. These
statements are proved over every field.
\end{abstract}

\section{Two obstructions read directly from the ternary operation}
Let a proper augmented nonpositive dg algebra $D$ over $k$ have
cohomology $k\oplus V\oplus W[1]$ with $V,W$ in degree zero before
shifting, finite dimensional, and zero products between all reduced
elements. A strictly unital augmented minimal model has only
$m_3:V^{\otimes3}\to W$ besides ordinary multiplication, by degrees.
The normalized Hochschild coboundaries into $W$ vanish on reduced
degree-zero triples. Write $X=V^*$ and
$\Phi:W^*\to X^{\otimes3}$ for the dual tensor, using a consistent
word convention. The reduced bar dual is
\[
 K=T(X\oplus W^*),\quad |X|=1,\quad |W^*|=2,\qquad
 dX=0,\quad d\eta=\Phi(\eta).
\]
Changing the overall sign on degree-two generators or reversing words
only changes conventions. The augmentation generates $D_{\fd}(D)$ by
truncation and finite length, so this category identifies with
$\per K$ by the functor $\RH_D(k,-)$.

\begin{prop}\label{obstruction}
If $\Phi$ has a nonzero kernel or if its image is contained in
$X_0^{\otimes3}$ for a proper subspace $X_0\subset X$, the profile
poset has exactly two elements. For a nonzero reduced cohomology
algebra the formal comparison likewise has exactly two.
\end{prop}
\begin{proof}
In the first case choose a basis containing a kernel vector $v$ in
degree two. No generator differential contains $v$ and $dv=0$. In
the second case choose $v$ in a complement to $X_0$ in degree one.
In either case $K=C*T(kv)$ as a dg algebra, so the outer derivation
killing $C$ and differentiating $v$ is a chain map. The augmentation
of $C$ retracts onto $T(kv)$, so $H^*K$ is infinite. By the two-cut
theorem $\RH_K(P,E)\simeq E\ten_KP^\vee$ is finite exactly when
$H^*E$ or $H^*P^\vee$ is finite. Accordingly every detector has the full
profile or $\UU_K$, with endpoints realized by $0,K$. The free module
$K$ as source distinguishes them. For a zero differential on the whole
tensor algebra any reduced generator works.
\end{proof}

The smallest possible $X_0$ is obtained by taking the span of all
one-slot outputs of contractions of the tensors in $\operatorname{im}\Phi$
against covectors in their other two slots. Indeed membership in
$X_0\otimes X\otimes X$ and its two permutations is equivalent, by
a basis adapted to $X_0$, to membership in $X_0^{\otimes3}$.
This is a field-linear test. The converse of the proposition is false:
$dy=x^3$ on one generator of each degree has injective tensor and
uses the whole space. The earlier proof for Example 3.11 gives exactly
two profiles in that case.

\section{A flat matrix family}
Fix $r,m\ge1$ and a matrix $C=(c_{ji})\in M_{m\times r}(k)$. Define
$D(C)$ with degree-zero basis $1,a,b_1,\ldots,b_r,h$ and degree-minus-one
basis $e,w_1,\ldots,w_m$. Together with the unit the entire nonzero
multiplication and differential table is
\[
 a^2=h,\qquad b_i e=e b_i=\sum_j c_{ji}w_j,\qquad d e=h .
\]
Triple radical products vanish by construction; $h$ annihilates the
radical, giving the Leibniz identity also on two negative factors.
Thus this is a graded-commutative dg algebra even in characteristic two.
Cohomology consists of the classes of $1,a,b_i,w_j$ with all reduced
products zero. The underlying complex splits off $k e\xrightarrow{\sim}kh$.
The same formulas over $k[c_{ji}]$ give a bounded complex of finite
free modules, with free constant graded cohomology under arbitrary base
change. Applying the contraction $h\mapsto e$ gives
$m_3(b_i,a,a)=\sum_j c_{ji}w_j$ and
$m_3(a,a,b_i)=-\sum_j c_{ji}w_j$ in the convention of the baseline.
Hence $C\ne0$ gives a nonzero Hochschild class in every characteristic.

\begin{thm}\label{rank}
For $D(C)$, the profile poset differs from that of formal cohomology
if and only if $r=m=\rank C$. In the other cases it has exactly two
elements. In the exceptional case the entire poset is the two-row
classification of \othernote{normal_localisations}{Bounded profiles and
finite-dimensional points of universal localisations}, with $g=2r$ and
the simultaneous exchange of the two generators in each pair.
In particular it contains an infinite antichain over every field.
\end{thm}
\begin{proof}
The cubic differential is $d y_j=[x^2,\sum_i c_{ji} z_i]$ on the dual
generators. Row and column basis changes put $C$ in rank normal form
over $k$. If its rank $d<m$, there is a free degree-two cycle $y_j$.
If $d<r$, there is a free degree-one cycle $z_i$. Proposition
\ref{obstruction} applies. When $r=m=d$, the filtered bar computation
in the normal-localisation note, with $n=2,Q=I$, gives the Koszul dual
$A_{2,r,I}$. Explicitly the leading words $xxz_i$ have no proper
overlaps, and in each fixed expanded word their replacements resolve
independent contractible pairs. This gives the quasi-isomorphism to
the quotient by the commutators. All hypotheses of the localisation
theorem are proved there, including perfect lifts on both sides.
There are infinitely many distinct simple orbits: let all free
generators act by the class of $T$ on $k[T]/p$ for a monic irreducible
$p$. These modules are simple, fixed by the exchange, and distinct
for distinct $p$. There are infinitely many choices even over finite
fields. Singleton supports are pairwise incomparable.
\end{proof}

In particular full rank of a rectangular matrix is insufficient. Nonzero
higher products of arbitrarily large rank occur with just two profiles.
In the square family the locus with an infinite antichain is exactly
the principal open $\det C\ne0$. This is a proved statement for this
family, not a universal semicontinuity principle.

\paragraph{Quantitative check and ordinary realisation.}
The normal words after rank reduction avoid $xxz_i$ for $1\le i\le d$,
with $r+1$ letters of degree one and $m-d$ additional degree-two letters.
The forbidden patterns have no overlaps. Appending a degree-one or
degree-two letter gives the recurrence with series
\[
 \sum_{n\ge0}\dim\Ext_{D(C)}^n(k,k)t^n
 =\big[1-(r+1)t-(m-d)t^2+d t^3\big]^{-1}.
\]
For the subtraction, any admissible prefix followed by $xxz_i$ has no
other occurrence straddling that suffix, so the subtracted copies
are disjoint. The filtered argument remains valid with the free
degree-two letters: one may take the free product of the active
quasi-isomorphism with the free cycles. Its alternating augmentation-ideal
tensor decomposition preserves quasi-isomorphisms over $k$.

All these kernels are strictly two-term and proper. The manuscript's
Theorem 3.7 gives full-profile ordinary one-map universal localisations
of dimensions $3(r+3)+(m+1)\to4(r+3)$. Its Corollary 3.10 realises every
profile above modulo perfects as a marked bounded defect, with ambient
rings allowed to vary with the detector. At $r=m=1$ this recovers the
six-dimensional result. For $r=m>1$ it separates the entire determinantal
divisor, including nonformal fibres, from its complement.

\paragraph{Status and comparison.}
The passage from relation coefficients to higher products has precedents
in \href{https://arxiv.org/abs/math/0606144}{Lu--Palmieri--Wu--Zhang};
the nonoverlap computation uses the mechanism of
\href{https://doi.org/10.1016/0021-8693(82)90104-1}{Anick}.
The conclusion here concerns eventual Hom for all finite dg modules
and the resulting marked defects. It does not classify arbitrary cubic
tensors of full variable span. The companion script reports its finite
ranges and compares independent strict and minimal bar calculations.
\end{document}
